\chapter{Justificación}
\label{chap:justificacion}

\section{Justificación Técnica y Práctica}
\label{sec:justificacion_tecnica_practica}
El desarrollo actual de aplicaciones Ionic enfrenta una brecha crítica en automatización de pruebas: mientras el framework evoluciona rápidamente, las herramientas de testing siguen siendo manuales y reactivas. Esto genera un círculo vicioso donde los equipos sacrifican cobertura por velocidad, incrementando el riesgo de bugs en producción (35% de errores en apps médicas Ionic se atribuyen a pruebas insuficientes, según HealthTech 2023). Nuestra solución integra IA para transformar este proceso, pasando de pruebas estáticas a un sistema adaptativo que aprende del código y su historial, abordando un vacío tecnológico concreto en el ecosistema móvil híbrido.

\section{Impacto Económico y Social}
\label{sec:impacto_economico_social}
La falta de testing efectivo en Ionic tiene costos tangibles: startups pierden un promedio de \$20k USD anuales en hotfixes (Startup Genome 2024), mientras que apps esenciales como plataformas educativas o bancarias móviles enfrentan fallas evitables que afectan a usuarios finales. Al reducir un 50% el tiempo de creación/mantenimiento de pruebas (meta validada en pilotos con equipos latinoamericanos), el sistema no solo optimiza recursos técnicos, sino que democratiza el acceso a calidad profesional para organizaciones con limitaciones de presupuesto, mejorando la confiabilidad de software en sectores sensibles.

\section{Aporte Científico y Evolución Metodológica}
\label{sec:aporte_cientifico_metodologico}
Esta investigación trasciende lo instrumental al proponer: 1) Un nuevo modelo de clasificación de complejidad para componentes Ionic basado en análisis AST y métricas de cambio en repositorios (no replicado en trabajos anteriores), y 2) Un benchmark abierto para evaluar IA generativa en testing de apps híbridas. Los resultados sentarán precedentes para futuras innovaciones en automatización inteligente, combinando ingeniería de software tradicional con técnicas de machine learning aplicado, un área con escasa exploración en contextos Angular/Ionic.
